\documentclass[10pt,a4paper,landscape]{article}

\usepackage[T1]{fontenc}
\usepackage[utf8]{inputenc}
\usepackage[spanish]{babel}
\usepackage{lmodern}
\usepackage{geometry}
\usepackage[table]{xcolor}
\usepackage{booktabs}
\usepackage{tabularx}
\usepackage{array}
\usepackage{graphicx}
\usepackage[hidelinks]{hyperref}

\geometry{top=1.25cm,bottom=1.15cm,left=1.45cm,right=1.45cm}
\setlength{\parindent}{0pt}
\setlength{\parskip}{0.35em}
\renewcommand{\arraystretch}{1.28}
\definecolor{upcblue}{HTML}{0076A8}
\definecolor{deepblue}{HTML}{174A6E}
\definecolor{lightblue}{HTML}{EAF4F8}
\newcolumntype{Y}{>{\raggedright\arraybackslash}X}
\newcolumntype{C}[1]{>{\centering\arraybackslash}p{#1}}

\hypersetup{
  pdftitle={Rúbrica 2 - Autoevaluación de la planificación temporal},
  pdfauthor={Víctor Abelló Torres}
}

\begin{document}
\thispagestyle{empty}
\begin{minipage}[c]{0.15\textwidth}
  \includegraphics[width=0.95\linewidth]{../../latex/assets/upc-logo.pdf}
\end{minipage}
\hfill
\begin{minipage}[c]{0.81\textwidth}
  {\color{deepblue}\LARGE\bfseries Rúbrica 2 · Autoevaluación}\\[0.15em]
  {\large Planificación temporal · Víctor Abelló Torres · 28/09/2026}
\end{minipage}

\vspace{0.65cm}
\begin{tabularx}{\textwidth}{>{\bfseries}p{4.1cm} C{2.3cm} Y}
\toprule
\rowcolor{deepblue}
\color{white}Criterio & \color{white}Valoración & \color{white}Justificación \\
\midrule
Descripción de las tareas & 4 · Bien alcanzado & Se definen 23 tareas con resultado de cierre, estimación individual en horas, predecesoras y recursos humanos y materiales. Las dependencias también se expresan como cadenas lógicas y se resumen en dos tablas. \\
\rowcolor{lightblue}
Estimaciones y Gantt & 4 · Bien alcanzado & El total es de 540 horas; las estimaciones técnicas son inferiores a 50 horas y se justifica la tarea transversal de memoria. El Gantt es coherente con las tablas y muestra concurrencia entre modelos, evaluación, documentación y seguimiento. \\
Gestión del riesgo: planes alternativos y obstáculos & 4 · Bien alcanzado & Cinco riesgos tienen señal de activación, tareas alternativas, impacto temporal y recursos. Los planes B no se incluyen en el Gantt base y se dispone de margen y de una reducción explícita del alcance opcional para proteger la fecha final. \\
\rowcolor{lightblue}
Organización & 4 · Bien alcanzado & Hay portada con los datos relevantes, cinco secciones enlazadas, tablas y figura numeradas, páginas numeradas y dos finales diferenciados: técnico y académico. \\
Claridad & 4 · Bien alcanzado & El documento comienza con los datos globales, define abreviaturas de recursos y emplea los mismos códigos T01--T23 en texto, tablas y Gantt, lo que facilita seguir la secuencia. \\
\rowcolor{lightblue}
Redacción & 4 · Bien alcanzado & Se ha realizado una revisión ortográfica y terminológica en castellano y se distingue entre fechas confirmadas, previsiones y datos pendientes de confirmación. \\
\midrule
\textbf{Autoevaluación global} & \textbf{4 / 4} & \textbf{El documento cubre todos los indicadores del nivel máximo de la rúbrica. La fecha oficial de lectura y el dispositivo edge permanecen correctamente marcados como pendientes.} \\
\bottomrule
\end{tabularx}

\vfill
\footnotesize
Esta autoevaluación acompaña al documento \emph{Entrega 2 de GEP: planificación temporal}. La calificación propuesta es una valoración del autor y queda sujeta a la evaluación docente.
\end{document}
